% ============================================================
%  Sistemas Operativos — Unidad 1, Clase 1
%  Introducción a los Sistemas Operativos
%  Docente: Rodrigo Fuchs Miranda
%  UTEC — Sede Rivera — Licenciatura en Ingeniería de Datos e IA
% ============================================================
\documentclass[xcolor=svgnames]{beamer}

% -----------------------------------------------------------
% Codificación e idioma
% -----------------------------------------------------------
\usepackage[utf8]{inputenc}
\usepackage[T1]{fontenc}
\usepackage[spanish]{babel}

% -----------------------------------------------------------
% Gráficos, colores y diagramas
% -----------------------------------------------------------
\usepackage{xcolor}
\usepackage{graphicx}
\usepackage{tikz}

% -----------------------------------------------------------
% Hyperlinks y ajustes de layout
% -----------------------------------------------------------
\usepackage{hyperref}
\hypersetup{colorlinks=true, linkcolor=blue}
\usepackage{adjustbox}
\usepackage{calc}

% -----------------------------------------------------------
% Símbolos
% -----------------------------------------------------------
\usepackage{pifont}
\usepackage{enumitem}
\setbeamertemplate{frametitle continuation}{\relax}
% -----------------------------------------------------------
% Tema y paleta de colores
% -----------------------------------------------------------
\usetheme{Madrid}

\definecolor{mqred}{RGB}{166, 25, 46}
\definecolor{mqdeepred}{RGB}{118, 35, 47}
\definecolor{mqgray}{RGB}{55, 58, 54}
\definecolor{mqlightgray}{RGB}{237, 235, 229}
\definecolor{mqmagenta}{RGB}{198, 0, 126}

\usecolortheme[named=mqred]{structure}
\setbeamercolor{title in head/foot}{bg=mqlightgray, fg=mqgray}
\setbeamercolor{author in head/foot}{bg=mqdeepred}
\setbeamercolor{page number in head/foot}{bg=mqdeepred, fg=mqlightgray}

% -----------------------------------------------------------
% Pie de página personalizado
% -----------------------------------------------------------
\makeatletter
\setbeamertemplate{footline}{
  \leavevmode%
  \hbox{%
    \begin{beamercolorbox}[wd=.5\paperwidth,ht=2.25ex,dp=1ex,center]{author in head/foot}%
      \usebeamerfont{author in head/foot}\insertshortauthor\expandafter\ifblank\expandafter{\beamer@shortinstitute}{}{~~(\insertshortinstitute)}
    \end{beamercolorbox}%
    \begin{beamercolorbox}[wd=.4\paperwidth,ht=2.25ex,dp=1ex,center]{title in head/foot}%
      \usebeamerfont{title in head/foot}\insertshorttitle
    \end{beamercolorbox}%
    \begin{beamercolorbox}[wd=.1\paperwidth,ht=2.25ex,dp=1ex,center]{page number in head/foot}%
      \usebeamerfont{page number in head/foot}\insertframenumber{} / \inserttotalframenumber
    \end{beamercolorbox}%
  }%
  \vskip0pt%
}
\makeatother
\beamertemplatenavigationsymbolsempty

% -----------------------------------------------------------
% Comandos para imágenes
% -----------------------------------------------------------
\newcommand{\imgfill}[1]{%
  \begin{center}
    \adjustbox{max height=0.5\textheight, keepaspectratio}{%
      \includegraphics{#1}%
    }%
  \end{center}
}
\newcommand{\imgfillAuto}[1]{%
  \smallskip
  \newlength{\available}
  \setlength{\available}{\dimexpr \textheight - \pagetotal - 1em \relax}
  \begin{center}
    \includegraphics[width=\textwidth, height=\available, keepaspectratio]{#1}%
  \end{center}
  \smallskip
}

% -----------------------------------------------------------
% Metadatos del documento
% -----------------------------------------------------------
\title[SO — Unidad 1]{\textbf{Sistemas Operativos} \\ \textit{Unidad 1, Clase 1: Introducción a los Sistemas Operativos}}
\author[Rodrigo Fuchs Miranda]{Rodrigo Fuchs Miranda}
\institute[Lic.\ en Ing.\ de Datos e IA]{Licenciatura en Ingeniería de Datos e Inteligencia Artificial}
\date{9 de marzo de 2026}

% ============================================================
\begin{document}
% ============================================================

% -----------------------------------------------------------
% Portada
% -----------------------------------------------------------
\begin{frame}
  \titlepage
  \begin{center}
  \vspace{-10pt}
    \includegraphics[height=3.5cm]{logo.jpg}
  \end{center}
\end{frame}

% -----------------------------------------------------------
% Índice
% -----------------------------------------------------------
\begin{frame}{¿Qué vemos hoy?}
  \tableofcontents
\end{frame}

% ============================================================
\section{Presentación de la asignatura}
% ============================================================

\begin{frame}{Presentación de la asignatura}
  \begin{columns}[T]
    \begin{column}{0.5\textwidth}
      \textbf{Clases:}
      \begin{itemize}
        \item[\textcolor{mqred}{\ding{228}}] Lunes: 10:10 -- 12:10
        \item[\textcolor{mqred}{\ding{228}}] Viernes: 10:10 -- 12:10
      \end{itemize}
      \vspace{0.4cm}
      \textbf{Modalidad:}
      \begin{itemize}
        \item[\textcolor{mqred}{\ding{72}}] Clase teórica: Salón 1A/1B
        \item[\textcolor{mqred}{\ding{72}}] Clase práctica:\\ Lunes: Lab. de Informática - UTU \\ Viernes: Lab. de Informática – Edificio Principal
      \end{itemize}
    \end{column}
    \begin{column}{0.5\textwidth}
      \textbf{Evaluación:}
      \begin{itemize}
        \item[\textcolor{mqred}{\ding{51}}] \textbf{PE} (30\,\%) — 11 de mayo
        \item[\textcolor{mqred}{\ding{51}}] \textbf{SE} (30\,\%) — 26 de junio
        \item[\textcolor{mqred}{\ding{51}}] \textbf{EC} (40\,\%) — Trabajos en clase y entregas a lo largo del semestre
      \end{itemize}
    \end{column}
  \end{columns}
\end{frame}

% -----------------------------------------------------------
\begin{frame}{Objetivos del curso}
  Al finalizar la asignatura, el estudiante será capaz de:
  \begin{itemize}
    \item[\textcolor{mqred}{\ding{228}}] \textbf{Comprender} la historia, evolución y funciones de los Sistemas Operativos.
    \item[\textcolor{mqred}{\ding{228}}] \textbf{Analizar} la administración de procesos, memoria y entrada/salida.
    \item[\textcolor{mqred}{\ding{228}}] \textbf{Conocer} los sistemas de archivos y los mecanismos de seguridad.
    \item[\textcolor{mqred}{\ding{228}}] \textbf{Explorar} técnicas de virtualización y contenedores.
  \end{itemize}
\end{frame}

% -----------------------------------------------------------
\begin{frame}{Contenidos y bibliografía}
  \begin{columns}[T]
    \begin{column}{0.55\textwidth}
      \begin{itemize}
        \item[\textcolor{mqred}{\ding{228}}] \textbf{Unidad 1:} Introducción a los S.O.
        \item[\textcolor{mqred}{\ding{228}}] \textbf{Unidad 2:} Administración de Procesos
        \item[\textcolor{mqred}{\ding{228}}] \textbf{Unidad 3:} Gestión de Memoria
        \item[\textcolor{mqred}{\ding{228}}] \textbf{Unidad 4:} Administración de E/S
        \item[\textcolor{mqred}{\ding{228}}] \textbf{Unidad 5:} Sistemas de Archivos y Seguridad
        \item[\textcolor{mqred}{\ding{228}}] \textbf{Unidad 6:} Virtualización y Contenedores
      \end{itemize}
    \end{column}
    \begin{column}{0.45\textwidth}
      \begin{center}
        \begin{tikzpicture}
          \node[draw=black, thick] {\includegraphics[width=0.45\linewidth]{biblio1.png}};
        \end{tikzpicture}
        \hspace{0.3cm}
        \begin{tikzpicture}
          \node[draw=black, thick] {\includegraphics[width=0.46\linewidth]{biblio2.png}};
        \end{tikzpicture}
      \end{center}
    \end{column}
  \end{columns}
\end{frame}

% ============================================================
\section{¿Para qué sirve un Sistema Operativo?}
% ============================================================

% -----------------------------------------------------------
%  Pregunta disparadora
% -----------------------------------------------------------
\begin{frame}{Una pregunta para arrancar}
  \begin{center}
    \Large
    \textbf{¿Qué pasaría si tu celular\\no tuviera sistema operativo?}
  \end{center}
  \vspace{0.5cm}
  \begin{itemize}
    \item No habría aplicaciones en ejecución.
    \item No habría pantalla táctil, Wi-Fi ni notificaciones.
    \item El hardware estaría ahí, pero \textbf{nadie lo gestionaría}.
  \end{itemize}
  \vspace{0.4cm}
  \begin{center}
    \textit{El S.O. es la capa que hace que el hardware sea útil.}
  \end{center}
\end{frame}

% -----------------------------------------------------------
%  El S.O. en el contexto de la máquina
% -----------------------------------------------------------
\begin{frame}{La computadora como un sistema en capas}
  \begin{columns}[T]
    \begin{column}{0.55\textwidth}
      Una computadora moderna integra:
      \begin{itemize}
        \item Procesador(es), memoria y dispositivos de E/S
        \item Tal como vimos en \textbf{Arquitectura de Computadoras}
      \end{itemize}
      \vspace{0.3cm}
      El \textbf{Sistema Operativo} es la capa intermedia entre el hardware y las aplicaciones:
      \begin{itemize}
        \item Abstrae los detalles del hardware
        \item Administra los recursos disponibles
        \item Brinda servicios a los programas de usuario
      \end{itemize}
    \end{column}
    \begin{column}{0.4\textwidth}
      \begin{center}
        \includegraphics[width=\columnwidth, keepaspectratio]{so_interfaz.png}
      \end{center}
      % Figura: diagrama de capas (hardware → SO → aplicaciones → usuario)
    \end{column}
  \end{columns}
\end{frame}

% -----------------------------------------------------------
%  Definición dual: máquina extendida + administrador
% -----------------------------------------------------------
\begin{frame}[allowframebreaks]{¿Qué es un Sistema Operativo?}

  El autor Tanenbaum define el S.O.\ desde \textbf{dos perspectivas complementarias}:

  \vspace{0.3cm}

  \begin{block}{1. El S.O. como máquina extendida}
    Oculta los detalles complejos del hardware y ofrece al programador
    herramientas más sencillas para trabajar.\\[4pt]
    \textit{Analogía:} así como un auto tiene volante y pedales (y no necesitás
    entender el motor para manejarlo), el S.O.\ te da \textbf{archivos,
    procesos y conexiones de red} sin que tengas que programar
    directamente sobre el disco, la CPU o la placa de red.
  \end{block}

  \vspace{0.3cm}

  \begin{block}{2. El S.O. como administrador de recursos}
    Gestiona y reparte \textbf{CPU, memoria y dispositivos de E/S}
    entre todos los programas que se ejecutan al mismo tiempo,
    garantizando que ninguno bloquee ni dañe al resto.
  \end{block}

\end{frame}

% -----------------------------------------------------------
\begin{frame}{¿Qué es un Sistema Operativo? }

  En síntesis: el S.O.\ es el \textbf{software} que gestiona los recursos del \textbf{hardware} y ofrece a los programas una interfaz estandarizada para utilizarlos.

  \vspace{0.3cm}

  \begin{center}
    \includegraphics[width=0.75\textwidth, keepaspectratio]{so_interfaz2.png}
  \end{center}
  % Figura: diagrama clásico de Tanenbaum mostrando el S.O. entre hardware y aplicaciones

\end{frame}

% -----------------------------------------------------------
%  Llamadas al sistema
% -----------------------------------------------------------
\begin{frame}{Las llamadas al sistema (\textit{system calls})}

  Los programas de usuario no pueden acceder directamente al hardware.\\
  Para pedir un servicio al S.O., utilizan las \textbf{llamadas al sistema}.

  \vspace{0.3cm}

  \begin{exampleblock}{Ejemplo cotidiano}
    Cuando una aplicación quiere guardar un archivo, no escribe directamente en el disco. Hace una \textbf{llamada al sistema} (\texttt{write}) y el S.O.\ se encarga del resto.
  \end{exampleblock}

  \vspace{0.3cm}

  \begin{itemize}
    \item Las llamadas al sistema son la \textbf{interfaz formal} entre los programas y el S.O.
    \item Veremos esto en detalle más adelante en el curso.
  \end{itemize}

\end{frame}

% ============================================================
\section{Interfaces de usuario y modos de operación}
% ============================================================

\begin{frame}{Interfaz de usuario: Shell vs.\ GUI}
  Los usuarios no interactúan directamente con el \textbf{kernel}, sino a través de programas de interfaz:

  \vspace{0.3cm}

  \begin{columns}[T]
    \begin{column}{0.5\textwidth}
      \begin{block}{Shell (interfaz de texto)}
        Intérprete de comandos: el usuario escribe instrucciones que el shell traduce y envía al S.O.\\
        \textit{Ej.: bash, zsh, PowerShell}
      \end{block}
    \end{column}
    \begin{column}{0.5\textwidth}
      \begin{block}{GUI (\textit{Graphical User Interface})}
        Entorno visual: ventanas, íconos y ratón.\\
        \textit{Ej.: Windows 11, macOS, GNOME}
      \end{block}
    \end{column}
  \end{columns}

  \vspace{0.4cm}
  \begin{alertblock}{Importante}
    El shell y la GUI \textbf{no son parte del kernel}: son programas de usuario que usan los servicios del S.O.
  \end{alertblock}
\end{frame}

% -----------------------------------------------------------
\begin{frame}{Modo kernel vs.\ modo usuario}

  Para proteger el sistema, la CPU opera en dos modos distintos:

  \vspace{0.3cm}

  \begin{columns}[T]
    \begin{column}{0.5\textwidth}
      \begin{block}{Modo kernel (supervisor)}
        \begin{itemize}
          \item El núcleo del S.O.\ tiene \textbf{acceso completo} al hardware.
          \item Puede ejecutar instrucciones privilegiadas.
          \item Gestiona memoria, dispositivos e interrupciones.
        \end{itemize}
      \end{block}
    \end{column}
    \begin{column}{0.5\textwidth}
      \begin{block}{Modo usuario}
        \begin{itemize}
          \item Las aplicaciones se ejecutan con \textbf{privilegios restringidos}.
          \item Si necesitan algo privilegiado, hacen una llamada al sistema.
          \item El S.O.\ cambia a modo kernel, atiende la solicitud y regresa.
        \end{itemize}
      \end{block}
    \end{column}
  \end{columns}

  \vspace{0.4cm}
  Esta separación evita que un programa dañe o bloquee toda la máquina.

\end{frame}

% ============================================================
\section{Clasificación de los Sistemas Operativos}
% ============================================================

\begin{frame}{¿De qué tipo puede ser un S.O.?}

  Los sistemas operativos se clasifican según su \textbf{propósito y entorno de uso}:

  \vspace{0.2cm}

  \begin{itemize}
    \item[\textcolor{mqred}{\ding{228}}] \textbf{Por lotes (\textit{batch}):} trabajan en secuencia sin interacción del usuario. \textit{Ej.: primeros mainframes.}
    \item[\textcolor{mqred}{\ding{228}}] \textbf{De tiempo compartido (\textit{timesharing}):} varios usuarios interactivos comparten la CPU. \textit{Ej.: UNIX, CTSS.}
    \item[\textcolor{mqred}{\ding{228}}] \textbf{De tiempo real:} responden en plazos muy acotados y predecibles. \textit{Ej.: control de robots, aviónica.}
    \item[\textcolor{mqred}{\ding{228}}] \textbf{Empotrados (\textit{embedded}):} dispositivos con recursos limitados y función específica. \textit{Ej.: electrodomésticos inteligentes.}
    \item[\textcolor{mqred}{\ding{228}}] \textbf{Distribuidos:} cooperan en red formando un sistema coordinado.
    \item[\textcolor{mqred}{\ding{228}}] \textbf{Móviles:} orientados a smartphones y tablets. \textit{Ej.: Android, iOS.}
  \end{itemize}

\end{frame}

% ============================================================
\section{Evolución histórica de los Sistemas Operativos}
% ============================================================

\begin{frame}{Evolución histórica: panorama general}
  \vspace{0.2cm}
  \setbeamercolor{gen1}{bg=mqred!15, fg=black}
  \setbeamercolor{gen2}{bg=mqred!25, fg=black}
  \setbeamercolor{gen3}{bg=mqred!38, fg=black}
  \setbeamercolor{gen4}{bg=mqred!52, fg=black}
  \setbeamercolor{gen5}{bg=mqred!70, fg=white}

  \begin{beamercolorbox}[wd=\textwidth, ht=0.55cm, dp=0.2cm, center, rounded=true]{gen1}
    \textbf{1ª gen. (1945--1955)} \quad Sin S.O., tubos de vacío
  \end{beamercolorbox}\vspace{0.15cm}

  \begin{beamercolorbox}[wd=\textwidth, ht=0.55cm, dp=0.2cm, center, rounded=true]{gen2}
    \textbf{2ª gen. (1955--1965)} \quad Transistores y sistemas en lote
  \end{beamercolorbox}\vspace{0.15cm}

  \begin{beamercolorbox}[wd=\textwidth, ht=0.55cm, dp=0.2cm, center, rounded=true]{gen3}
    \textbf{3ª gen. (1965--1980)} \quad Multiprogramación y timesharing
  \end{beamercolorbox}\vspace{0.15cm}

  \begin{beamercolorbox}[wd=\textwidth, ht=0.55cm, dp=0.2cm, center, rounded=true]{gen4}
    \textbf{4ª gen. (1980--2000)} \quad PC y GUI
  \end{beamercolorbox}\vspace{0.15cm}

  \begin{beamercolorbox}[wd=\textwidth, ht=0.55cm, dp=0.2cm, center, rounded=true]{gen5}
    \textbf{5ª gen. (2000--hoy)} \quad Movilidad y computación en la nube
  \end{beamercolorbox}
\end{frame}

% -----------------------------------------------------------
\subsection{Primera generación (1945--1955)}
\begin{frame}{1ª generación (1945--1955): las máquinas sin S.O.}
  \begin{itemize}
    \item Computadoras construidas con \textbf{tubos de vacío}. Sin sistema operativo.
    \item El usuario era también el programador: conectaba cables y configuraba la máquina a mano (plugboards) o ingresaba código binario directamente.
    \item Un solo usuario a la vez: \textbf{sin multitarea, sin memoria protegida}.
    \item Uso casi exclusivo en cálculos científicos y militares.
    \item \textbf{ENIAC (1946):} $\sim$18.000 válvulas, programación totalmente manual.
  \end{itemize}
  \vspace{0.3cm}
  \begin{center}
    \textit{El tiempo de CPU era tan valioso que el programador debía reservar la sala de máquinas con anticipación.}
  \end{center}
\end{frame}

% -----------------------------------------------------------
\subsection{Segunda generación (1955--1965)}
\begin{frame}{2ª generación (1955--1965): sistemas en lote}
  \begin{itemize}
    \item Reemplazo de tubos por \textbf{transistores}: máquinas más rápidas y confiables.
    \item Los trabajos se organizan en \textbf{lotes de tarjetas perforadas}, minimizando el tiempo ocioso de la CPU.
    \item Aparecen los primeros \textbf{monitores residentes}: programas básicos que cargan automáticamente el siguiente trabajo del lote.
    \item No había interacción del usuario durante la ejecución.
    \item Ejemplo: \textbf{IBM 7094} con \textbf{FMS} (\textit{Fortran Monitor System}).
  \end{itemize}
  \begin{figure}
    \centering
    \includegraphics[width=\linewidth]{geracao2.png}
  \end{figure}
\end{frame}

% -----------------------------------------------------------
\begin{frame}{Estructura de un trabajo en FMS}
  Los trabajos se definían mediante \textbf{tarjetas de control}:
  \begin{itemize}
    \item \textbf{\$JOB} — inicio del trabajo, identifica al usuario
    \item \textbf{\$FORTRAN} — indica al sistema que compile el programa en Fortran
    \item \textbf{\$LOAD} — carga el programa compilado en memoria
    \item \textbf{\$RUN} — ejecuta el programa
    \item \textbf{\$END} — fin del trabajo; el monitor carga el siguiente lote
  \end{itemize}
  \begin{center}
    \includegraphics[width=0.55\textwidth]{fortran.png}
  \end{center}
  % Figura: diagrama de tarjetas de control del FMS (Tanenbaum, Fig. 1-3)
\end{frame}

% -----------------------------------------------------------
\subsection{Tercera generación (1965--1980)}
\begin{frame}{3ª generación (1965--1980): multiprogramación y timesharing}
  \begin{columns}[T]
    \begin{column}{0.65\textwidth}
    \small   
      \begin{itemize}
        \item \textbf{Multiprogramación:} varios programas cargados en memoria simultáneamente; mientras uno espera E/S, otro usa la CPU.
        \item \textbf{Timesharing:} múltiples usuarios interactivos comparten la CPU en turnos muy rápidos, simulando exclusividad.
        \item \textbf{IBM 360:} familia unificada para uso científico y comercial. Primer intento serio de estandarización.
        \item \textbf{MULTICS:} proyecto muy ambicioso, influyó en el diseño de sistemas posteriores.
        \item \textbf{UNIX (Bell Labs):} más simple y portable que MULTICS, se populariza en universidades y es ancestro de Linux, macOS y Android.
      \end{itemize}
    \end{column}
    \begin{column}{0.35\textwidth}
      \begin{center}
        \includegraphics[width=\columnwidth, keepaspectratio]{tercera.png}
      \end{center}
      % Figura: IBM 360 o diagrama de multiprogramación
    \end{column}
  \end{columns}
\end{frame}

% -----------------------------------------------------------
\subsection{Cuarta generación (1980--2000)}
\begin{frame}{4ª generación (1980--2000): la PC y la GUI}
  \begin{itemize}
    \item Nacen las \textbf{computadoras personales} (PC) asequibles para el público masivo, gracias a microprocesadores como el Intel 8080 y el Motorola 68000.
    \item Sistemas operativos de línea de comandos: \textbf{CP/M}, \textbf{MS-DOS}.
    \item Apple Macintosh (1984) y luego \textbf{Windows} popularizan la \textbf{GUI}: ventanas, íconos, ratón — sin necesidad de memorizar comandos.
    \item Amplia adopción en hogares, oficinas y educación.
    \item Aparece \textbf{Linux} (1991, Linus Torvalds): kernel libre que hoy sustenta servidores, Android y supercomputadoras.
  \end{itemize}
\end{frame}

% -----------------------------------------------------------
\subsection{Quinta generación (2000--presente)}
\begin{frame}{5ª generación (2000--hoy): movilidad y cloud}
  \begin{itemize}
    \item \textbf{Dispositivos móviles:} Android e iOS gestionan pantallas táctiles, GPS, cámara, conectividad y batería — todo en un bolsillo.
    \item \textbf{Computación en la nube:} los S.O.\ deben gestionar miles de servidores virtuales en grandes centros de datos.
    \item \textbf{Virtualización y contenedores:} Docker, Kubernetes y las máquinas virtuales permiten aislar y escalar servicios de forma eficiente.
    \item Linux y Windows Server dominan como base en infraestructura cloud.
    \item Los S.O.\ modernos priorizan \textbf{seguridad, escalabilidad y eficiencia energética}.
  \end{itemize}
\end{frame}

% ============================================================
\section{Actividad y cierre}
% ============================================================

\begin{frame}{Actividad: ¿Con o sin S.O.?}

  \textbf{De forma individual:} para cada situación, decidí si hay un S.O.\ involucrado o no, y anotá por qué.

  \vspace{0.3cm}

  \begin{enumerate}
    \item Apretás el botón de encendido de tu celular.
    \item Una calculadora de bolsillo suma dos números.
    \item Netflix te recomienda una película.
    \item Un semáforo cambia de rojo a verde.
    \item Escribís un documento en Word.
  \end{enumerate}

  \vspace{0.3cm}
  \begin{alertblock}{Para pensar}
    ¿Qué criterio usaste para decidir si hay S.O.\ o no?
  \end{alertblock}

\end{frame}

\begin{frame}{Actividad: respuestas y discusión}

  \begin{itemize}
    \item[\textcolor{mqred}{\ding{51}}] \textbf{Celular encendiendo} — \textit{Sí.} Es el proceso de arranque (\textit{boot}) del S.O.
    \item[\textcolor{mqred}{\ding{55}}] \textbf{Calculadora de bolsillo} — \textit{No.} Es un circuito lógico dedicado, sin gestión de recursos ni procesos.
    \item[\textcolor{mqred}{\ding{51}}] \textbf{Netflix} — \textit{Sí.} Los servidores corren Linux; tu dispositivo usa su propio S.O.\ para mostrarte el video.
    \item[\textcolor{mqred}{\ding{228}}] \textbf{Semáforo} — \textit{Depende.} Los semáforos modernos e interconectados tienen un S.O.\ embebido de tiempo real. Los más simples son solo un temporizador sin S.O.
    \item[\textcolor{mqred}{\ding{51}}] \textbf{Word} — \textit{Sí.} Word es una aplicación que corre sobre el S.O., el cual gestiona memoria, teclado y disco.
  \end{itemize}

  \vspace{0.2cm}
  \begin{exampleblock}{Conclusión}
    No todo dispositivo electrónico tiene S.O. Lo que define su presencia es la necesidad de \textbf{gestionar recursos} y ejecutar \textbf{múltiples procesos}.
  \end{exampleblock}

\end{frame}

% -----------------------------------------------------------
\begin{frame}{Cierre y síntesis de la clase}

  \textbf{¿Qué vimos hoy?}

  \vspace{0.3cm}

  \begin{itemize}
    \item[\textcolor{mqred}{\ding{51}}] El S.O.\ actúa como \textbf{máquina extendida} (oculta la complejidad del hardware) y como \textbf{administrador de recursos} (gestiona CPU, memoria y E/S).
    \item[\textcolor{mqred}{\ding{51}}] Los programas se comunican con el S.O.\ a través de \textbf{llamadas al sistema}.
    \item[\textcolor{mqred}{\ding{51}}] La separación entre \textbf{modo kernel} y \textbf{modo usuario} protege al sistema.
    \item[\textcolor{mqred}{\ding{51}}] Los S.O.\ se clasifican según su entorno de uso: lote, tiempo compartido, tiempo real, embebido, distribuido, móvil.
    \item[\textcolor{mqred}{\ding{51}}] La evolución histórica va desde máquinas sin S.O.\ hasta entornos virtualizados en la nube.
  \end{itemize}

  \vspace{0.3cm}
  \begin{center}
    \textbf{Próxima clase:} funciones del S.O.\ en detalle y discusión guiada.
  \end{center}

\end{frame}

% ============================================================
\section{Glosario de términos}
% ============================================================

\begin{frame}{Glosario (I)}
  \begin{description}[labelwidth=3.2cm, leftmargin=!, font=\color{mqred}\bfseries]
    \item[Sistema Operativo] Software que gestiona el hardware de una computadora
      y provee servicios a los programas de usuario.
    \item[Kernel] Núcleo del S.O.; la parte que corre con máximos privilegios
      y controla directamente el hardware.
    \item[Hardware] Componentes físicos de la computadora: procesador, memoria,
      disco, teclado, pantalla, etc.
    \item[Proceso] Un programa en ejecución. El S.O.\ administra cuándo y cómo
      cada proceso usa la CPU y la memoria.
    \item[Abstracción] Simplificación de algo complejo: el S.O.\ oculta los
      detalles del hardware y ofrece herramientas más fáciles de usar.
    \item[Llamada al sistema] Mecanismo formal por el cual un programa le pide
      un servicio al S.O.\ (ej.: leer un archivo, enviar datos por la red).
  \end{description}
\end{frame}

\begin{frame}{Glosario (II)}
  \begin{description}[labelwidth=3.2cm, leftmargin=!, font=\color{mqred}\bfseries]
    \item[Modo kernel] Modo de operación de la CPU en el que el S.O.\ tiene
      acceso total al hardware. Solo el kernel puede ejecutar instrucciones
      privilegiadas.
    \item[Modo usuario] Modo restringido en el que corren las aplicaciones.
      Para acceder al hardware deben pasar por el kernel.
    \item[Shell] Programa de interfaz de texto que interpreta comandos del
      usuario y los envía al S.O.\ (ej.: bash, zsh).
    \item[GUI] \textit{Graphical User Interface}: interfaz visual con ventanas,
      íconos y ratón. No forma parte del kernel.
    \item[Multiprogramación] Técnica que permite tener varios programas cargados
      en memoria al mismo tiempo, aprovechando la CPU cuando uno espera E/S.
    \item[Timesharing] Variante de multiprogramación donde varios usuarios
      interactivos comparten la CPU en turnos muy rápidos.
  \end{description}
\end{frame}

\begin{frame}{Glosario (III)}
\small
  \begin{description}[labelwidth=3.2cm, leftmargin=!, font=\color{mqred}\bfseries]
    \item[E/S (Entrada/Salida)] Operaciones que intercambian datos entre la CPU
      y dispositivos externos: disco, teclado, pantalla, red, etc.
    \item[Sistema en lote] S.O.\ que procesa trabajos en secuencia, sin
      interacción del usuario durante la ejecución.
    \item[S.O.\ embebido] S.O.\ diseñado para dispositivos con recursos
      limitados y función específica (ej.: un electrodoméstico inteligente).
    \item[Virtualización] Técnica que permite ejecutar múltiples sistemas
      operativos o entornos aislados sobre el mismo hardware físico.
    \item[Contenedor] Entorno de ejecución aislado que comparte el mismo kernel
      del S.O.\ anfitrión. Más liviano que una máquina virtual.
      \textit{Ej.: Docker.}
    \item[UNIX] Sistema operativo creado en Bell Labs (1969), ancestro de Linux,
      macOS y Android. Base conceptual de gran parte del software moderno.
  \end{description}
\end{frame}

% ============================================================
\end{document}
% ============================================================
